\section{Introduction}

Formal verification of Multi-Agent Systems (MAS) requires logical formalisms capable of capturing how multiple agents interact, compete, and cooperate over time \cite{shoham2008multiagent,wooldridge2009introduction,DBLP:journals/tocl/MogaveroMPV14}. In this setting, strategic logics such as Alternating-time Temporal Logic (ATL$^\ast$) \cite{alur2002alternating} provide a natural framework for specifying what coalitions of agents can enforce, regardless of the behavior of others. This makes the logic particularly suitable for safety-critical and adversarial scenarios, where guarantees must hold even in the presence of uncontrollable or competing entities. Yet, writing correct ATL$^\ast$ specifications requires substantial expertise in formal methods, temporal logic, and game-theoretic semantics, making the construction of strategic specifications difficult to non-expert practitioners. This remains true even for ATL, the vanilla fragment of ATL$^\ast$. A major source of difficulty is that requirements are typically expressed in natural language, which is inherently ambiguous, both because its interpretation depends on context and because the same surface form may admit different structural readings \cite{saussure1916course,small2013lexical}. In formal verification, however, ambiguity cannot be tolerated: verification pipelines require well-formed formulas with precise syntax and unambiguous semantics. More generally, the gap between linguistic expression and logical form has long motivated the view that the surface grammar of a sentence may fail to make its underlying formal structure explicit \cite{frege1892sense,russell1905denoting,Russell1919-RUSTPO-65}. In the context of strategic reasoning, this mismatch becomes especially critical, since ambiguities may affect coalition attribution, temporal scope, and the intended interpretation of alternative strategic outcomes \cite{kulkarni2025dynamic}. These issues are further amplified when natural-language requirements are processed by Large Language Models (LLMs). Besides inheriting the intrinsic ambiguity of natural language, LLM inputs may also be under-specified or contain irrelevant task-related details, both of which can lead to multiple plausible formalizations of the same request. This is particularly problematic in the strategic setting, where even small differences in interpretation may produce logically distinct ATL or ATL$^\ast$ formulas with substantially different verification outcomes.

Recent advances in LLMs have generated growing interest in the automatic translation of natural-language requirements into formal specifications. Existing work, however, has focused predominantly on single-agent or purely temporal formalisms, such as Linear Temporal Logic (LTL) \cite{pnueli1977temporal}, Signal Temporal Logic (STL) \cite{maler2004monitoring}, and Computation Tree Logic (CTL) \cite{clarke1981design}. By contrast, the strategic setting of MAS remains largely unexplored. To the best of our knowledge, there is currently no established trained LLM, nor any publicly available dataset, specifically designed to support the translation of natural-language requirements into ATL or ATL$^\ast$ specifications.

\paragraph{Our Contributions.}
We introduce an open-weight, open-source LLM framework for translating natural-language requirements involving strategic reasoning into well-formed ATL/ATL$^\ast$ formulas. Our contributions are:

\begin{itemize}
    \item \textbf{First expert-authored NL-to-ATL/ATL$^\ast$ dataset.}
    We release, to the best of our knowledge, the first benchmark for translating natural-language requirements into ATL/ATL$^\ast$ formulae: 1100 instances hand-written by three domain experts and covering coalitional ability, adversarial outcomes, temporal objectives, and five ambiguity families: right dislocation, left dislocation, VP ellipsis, Right Node Raising, and quantifier scope. The dataset is approximately balanced across ATL/ATL$^\ast$ temporal operators and ambiguity types, and uses diverse realizations to avoid surface-template bias. 

    \item \textbf{Hybrid open-weight translation framework.}
    We propose an LLM-based NL-to-ATL/ATL$^\ast$ translation framework supporting both Azure OpenAI inference and private local deployment with open-weight models, enabling cloud-scale experimentation and privacy-sensitive industrial use.

    \item \textbf{Fine-tuning and judge reliability.}
    We fine-tune open-weight models on our dataset and compare them with strong few-shot proprietary baselines. Under the judge most aligned with human experts, small in-domain fine-tuned open-weight models match few-shot proprietary systems while improving transparency and deployment control. We also find that judge reliability runs \emph{opposite} to generator strength: the open-weight Llama-3.3-70B best tracks human verdicts, whereas the strongest proprietary models judge worst by over-rejecting input-grounded paraphrases, including their own outputs.

    \item \textbf{Verification-oriented integration.}
    We integrate the framework into \textsc{VITAMIN}~\cite{DBLP:conf/ifaamas/0001M25}, allowing natural-language strategic requirements to be translated into ATL/ATL$^\ast$ formulae and immediately used in an existing model-checking workflow.
    
\end{itemize}
We release the dataset, models, prompts, evaluation scripts, and orchestration code under an open-source license, providing a reusable foundation for LLM-assisted formal specification of Multi-Agent Systems.

\paragraph{Related works.} 
The problem of translating natural language specifications into temporal logics has received increasing attention in recent years \cite{cosler2023iterative,Finkbeiner2020TeachingTL,nayak2022req2spec,wu2022autoformalization}, particularly in the context of Linear Temporal Logic (LTL), Signal Temporal Logic (STL), and Computation Tree Logic (CTL). To the best of our knowledge, the most relevant peer-reviewed contributions along this research line are ~\cite{cosler2023nl2spec,fuggitti2023nl2ltl,DBLP:conf/icse/HeBNIG22,DBLP:journals/jss/ZrelliMAMSN26,mendoza2024translating}, which represent the current state of the art in mapping natural-language requirements to temporal logic specifications. The Fuggitti's and Cosler's works \cite{fuggitti2023nl2ltl,cosler2023nl2spec} primarily focus on LTL and explore different degrees of natural language abstraction and compositionality, while He et al. \cite{DBLP:conf/icse/HeBNIG22} and Zrelli et el. \cite{DBLP:journals/jss/ZrelliMAMSN26} respectively focus on STL and CTL. To the best of our knowledge, no prior peer-reviewed work has systematically addressed the problem of translating unrestricted natural language specifications into strategic temporal logics that explicitly capture agent coalitions and their strategic interaction. In this sense, our work constitutes the first dataset-driven and tool-supported study targeting natural language translation into ATL/ATL$^\ast$, thus extending existing approaches beyond purely temporal reasoning toward multi-agent strategic reasoning. Against Fugitti and Chakraborti work \cite{fuggitti2023nl2ltl}, we contend that the proposed natural language specifications remain overly aligned with the syntax and semantics of the target translation logic (LTL in that case). By contrast, our approach adopts natural language specifications with a higher expressive grade, allowing a wider range of semantic synonyms and syntactic realizations (e.g. left and right peripheral displacements, right node raising, verbal phrase ellipsis and quantifier scope ambiguities), adding complexity to the previously cited researches. Furthermore, while we acknowledge that the proposed translation mechanism can be valuable for progressively validating translation correctness in a compositional fashion, it arguably presupposes users who are already familiar with recursive construction over well-formed formulas and the associated composition rules. Therefore, our tool targets non-expert users who lack prior training in formal-methods notation and reasoning. The synthesis of correct LTL formulas from natural language specifications benefits our work, as the already cited literature in this area provides well established mappings for temporal operators. However, LTL lacks the syntactic and semantic machinery required to capture strategic behaviour, which is expressible in ATL/ATL$^\ast$. Accordingly, a central challenge of this work is to enable robust translation of such strategic operators, and their diverse natural language renderings.

%\paragraph{Outline.}Section~2 introduces preliminaries. Section~3 presents the dataset. Section~4 describes the proposed translation framework. Section~5 reports on evaluation and applications. Section~6 reports on the experiments. Section~7 concludes the work.